\section{数据集生成流水线}

\subsection{Instruction Data 生成}

讲者介绍了一个典型的 instruction data 生成流水线：

\begin{enumerate}
    \item \textbf{Seed Data}：起点可以是原始文本、已有的 QA 对、或用户真实查询
    \item \textbf{指令生成}：若只有原始文本（即答案），可通过 \textbf{back translation}（反向翻译）让 LLM 生成对应的问题
    \item \textbf{答案生成}：使用 LLM 基于指令生成回答
    \item \textbf{评分与过滤}：
    \begin{itemize}
        \item 使用启发式规则或 LLM-as-Judge 评估质量
        \item 数据去重（deduplication）
        \item 去污染（decontamination）——确保不包含测试集样本
        \item 关键词过滤——移除不需要的内容
    \end{itemize}
\end{enumerate}

\subsubsection{Instruction Following 数据的特殊处理}

Instruction following 是一类专注于\textbf{约束遵循}的数据。例如："用全小写英文回答"。生成此类数据时，可以将约束条件附加到 prompt 之后，让 LLM 生成回答，然后用\textbf{程序化测试}验证约束是否被满足（如语言检测、大小写检查），筛除不符合要求的样本。

\subsection{Preference Data 生成：Ultra Feedback 范式}

Preference data 的生成采用不同策略：

\begin{enumerate}
    \item 给定 prompt，查询\textbf{多个不同 LLM}（而非单一模型）
    \item 使用 Judge LLM 对每个回答评分
    \item 取最高分的回答作为 \textbf{chosen}，最低分的作为 \textbf{rejected}
    \item 去重、移除过短回答等后处理
\end{enumerate}

\begin{knowledgebox}{Decontamination（去污染）}
去污染是数据生成流水线中至关重要的一步。例如，IFEval 是一个广泛使用的 instruction following 评测集。在生成训练数据时，必须确保生成的样本与测试集不存在高相似度，否则会导致评测分数虚高但实际能力并未提升。
\end{knowledgebox}

\subsection{案例：Open Perfect Blend 数据集}

讲者展示了自己制作的 Open Perfect Blend 数据集，这是一个面向通用 fine-tuning 的混合数据集。其类别配比大致为：

\begin{itemize}
    \item \textbf{Math}（数学）：较大比例
    \item \textbf{Code}（代码）：较大比例
    \item \textbf{Chat}（通用对话）：中等比例
    \item \textbf{Instruction Following}：较小比例
\end{itemize}

这种混合比例（data mixture）的设计直接决定了模型在不同能力维度上的表现。

\subsection{Chat Template}

数据准备完成后的最后一步是应用 \textbf{Chat Template}。这是让模型理解"谁在说话"的关键结构：

\begin{lstlisting}[caption={Chat Template 示例（ChatML 格式）}]
<|im_start|>system
You are a helpful assistant.
<|im_end|>
<|im_start|>user
What is the capital of France?
<|im_end|>
<|im_start|>assistant
The capital of France is Paris.
<|im_end|>
\end{lstlisting}

特殊 token \texttt{<|im\_start|>} 和 \texttt{<|im\_end|>} 标记每条消息的起止，角色标签（system / user / assistant）明确说话者身份。多轮对话通过重复 user-assistant 对实现。

\begin{importantbox}{Chat Template 是 SFT 的核心学习目标}
Chat Template 的结构是 SFT 阶段模型需要学习的核心内容之一。正是这种结构化的对话格式让 base model 从"只会续写文本"进化为"能进行多轮对话的 assistant"。不同模型家族（Llama, Mistral, Qwen 等）使用不同的 chat template，但原理相同。
\end{importantbox}

\subsection{本章小结}

数据集生成是一个系统化的流水线，包括种子数据准备、指令/答案生成、评分过滤、去重去污染等步骤。Preference data 通过多模型采样 + Judge 评分的方式获取。最终需将数据格式化为 Chat Template 供 SFT 训练使用。

\newpage
